% !TeX program = xelatex
\documentclass[12pt,a4paper]{article}

\usepackage{fontspec}
\usepackage{polyglossia}
\setmainlanguage{russian}
\setotherlanguage{english}
\setmainfont{Times New Roman}
\newfontfamily\cyrillicfont{Times New Roman}
\setmonofont{Menlo}
\newfontfamily\cyrillicfonttt{Menlo}
\usepackage{geometry}
\usepackage{amsmath}
\usepackage{amssymb}
\usepackage{booktabs}
\usepackage{longtable}
\usepackage{array}
\usepackage{hyperref}

\geometry{margin=2.5cm}

\title{Методология формирования доменных коллекций с разметкой полезности и надежности документов}
\author{Григорьев Д. А., Чернышев Д. И.}
\date{09.09.2026}

\begin{document}
\sloppy

\maketitle

\begin{abstract}

В работе рассматривается методология формирования доменных коллекций для больших языковых моделей. 

\end{abstract}

\section{Введение}

Заполнится после.

\section{Литературный обзор}

Качество данных для предварительного обучения больших языковых моделей является критическим фактором, 
определяющим их итоговую производительность. 
За последние годы был предложен ряд крупных открытых корпусов, различающихся по источникам, объёму и методам очистки. 
В данном обзоре мы рассматриваем эволюцию подходов к созданию тренировочных датасетов, чтобы обосновать необходимость 
разработки инструментов для детальной разметки полезности и надёжности документов относительно предметных доменов.

Ранние работы, такие как \textbf{The Pile} \cite{pile} предложили подход, 
основанный на агрегации множества разнородных источников.
Он представляет собой англоязычный корпус объёмом около 825~ГБ, 
объединяющий научные статьи, книги, программный код, судебные документы и веб-страницы. 
Разделение данных в нём определяется источником происхождения, 
а не содержательными характеристиками отдельных документов.
\textbf{OSCAR} \cite{oscar} сделал шаг в сторону автоматической фильтрации веб-контента из Common Crawl
и ориентирован на построение крупных многоязычных корпусов 
с автоматическим определением языка и базовой фильтрацией веб-текстов. 
В более поздних версиях обработка была перенесена с отдельных строк на целые документы, 
а корпус дополнен метаданными о качестве и потенциально нежелательном содержании.

\textbf{RefinedWeb} \cite{refinedweb} демонстрирует, что тщательно очищенные веб-данные могут 
быть конкурентоспособны по отношению к корпусам, составленным из множества специализированных источников. 
Его обработка включает извлечение основного текста, языковую фильтрацию и удаление низкокачественных документов.
\textbf{CulturaX} \cite{culturax} развивает подходы к масштабируемой очистке многоязычных веб-корпусов.
Набор данных построен на основе mC4 и OSCAR и содержит более 6 трлн токенов для 167 языков. 
Его обработка включает языковую идентификацию, фильтрацию URL, очистку на основе текстовых метрик и дедупликацию, 
а пороги фильтрации определяются отдельно для каждого языка. 

\textbf{FineWeb} \cite{fineweb} продолжает развитие подходов к созданию крупных корпусов на основе Common Crawl.
Набор данных содержит около 15 триллионов токенов, полученных из Common Crawl.
Авторы подробно исследуют влияние отдельных этапов обработки, включая дедупликацию,
фильтрацию низкокачественных документов и удаление нежелательного содержимого.
На основе FineWeb также была сформирована коллекция \textbf{FineWeb-Edu},
содержащая тексты с высокой образовательной ценностью.
Для её построения использовалась модельная оценка качества документов,
что позволило выделить около 1,3 трлн токенов образовательного содержания.

\textbf{FineWeb2} \cite{fineweb2} адаптирует предложенный в FineWeb процесс обработки
для многоязычных веб-данных.
Полученный корпус занимает около 20~ТБ, содержит примерно 5 млрд документов и охватывает более 1000 языков.
FineWeb2 представляет собой один из наиболее масштабных открытых многоязычных корпусов,
однако его разметка также в первую очередь ориентирована на общее качество документов,
а не на их принадлежность конкретным предметным областям.

Рассмотренные работы демонстрируют постепенный переход от объединения заранее выбранных источников
к масштабируемой автоматической очистке и модельной оценке качества веб-документов.
При этом содержательная принадлежность документа к предметной области, его образовательная ценность,
надёжность и потенциальные риски обычно не представлены в виде единого набора признаков.
Это ограничивает возможности формирования специализированных обучающих коллекций
с заданными требованиями к тематике и качеству содержимого.


\section{Методология формирования доменных коллекций}

\subsection{Общая схема обработки}

Предлагаемый подход представляет собой последовательный пайплайн разметки и отбора документов из крупного веб-корпуса. 
На первом этапе документы группируются по доменам, после чего из корпуса формируется ограниченная выборка для первичной разметки. 
Каждый текст из этой выборки передается крупной языковой модели, которая присваивает ему набор тематических тегов в свободной форме.

Поскольку исходный набор тегов содержит большое количество близких и дублирующих формулировок, полученные теги дополнительно группируются в семантически однородные кластеры. 
Для каждого кластера выбирается репрезентативный тег, благодаря чему формируется фиксированное множество меток. 
На основе данной разметки обучается модель-теггер, которая затем применяется ко всему исходному корпусу.

Полученные предсказания используются для формирования коллекций, относящихся к заданным предметным областям. 
Кроме тематических тегов, каждому документу могут быть присвоены дополнительные характеристики, 
в том числе оценка образовательной ценности, токсичности и другие признаки, полезные при отборе данных для обучения языковых моделей.

\subsection{Обучение теггера}

Тематическая разметка веб-корпуса может быть получена непосредственно с помощью крупной языковой модели. 
Однако применение такой модели ко всему корпусу требует значительных вычислительных затрат. 
Поэтому в работе используется двухэтапная схема типа \textit{LLM--BERT}: 
крупная языковая модель формирует исходную разметку для ограниченной выборки документов, 
после чего на полученных данных обучается компактный тематический теггер, предназначенный для обработки полного корпуса.

Пусть
\begin{equation}
\mathcal{D} = {x_1, x_2, \ldots, x_N}
\end{equation}
--- выборка документов, используемая для построения обучающей разметки. 
Для каждого документа $x_n$ языковая модель $M$ формирует фиксированное число кратких тематических тегов:

\begin{equation}
M(x_n; q) ={{t}_{n1}, {t}_{n2}, \ldots, {t}_{nm}} = {T}(x_n),
\end{equation}

где $q$ --- инструкция для модели, а ${t}_{nj}$ --- свободный текстовый тег, описывающий одну из тем документа. 
Использование свободных тегов позволяет не задавать тематическую таксономию заранее и учитывать содержание конкретного корпуса. 
Вместе с тем такая разметка не является полностью согласованной: 
одна и та же тема может быть обозначена синонимичными формулировками, различными тегами. 

Перед построением фиксированного пространства классов выполняется нормализация тегов. 
Нормализация включает приведение к единому регистру, удаление лишних пробелов и специальных символов, 
унификацию разделителей и исключение дубликатов внутри одного документа.
После нормализации формируется множество уникальных свободных тегов $\mathcal{U}$.

Размер множества $\mathcal{U}$ может быть существенно больше числа тематических классов, 
необходимого для обучения классификатора. Например, теги, обозначающие одну предметную область, 
могут отличаться порядком слов, используемой терминологией или уровнем обобщения. 
Для преобразования открытого пространства тегов в конечную таксономию применяется их семантическая кластеризация.

Каждому нормализованному тегу $u \in \mathcal{U}$ сопоставляется векторное представление, 
полученное с помощью модели текстовых эмбеддингов $E$:

\begin{equation}
\mathbf{e}_u = E(u)
\end{equation}

Семантическая близость двух тегов определяется косинусным сходством:

\begin{equation}
s(u_i,u_j)
=
\cos\left(\mathbf{e}_{u_i},\mathbf{e}_{u_j}\right)
=
\mathbf{e}_{u_i}^{\mathsf{T}}\mathbf{e}_{u_j}.
\end{equation}

Соответствующее расстояние задаётся как

\begin{equation}
d(u_i,u_j)
=
1-s(u_i,u_j).
\end{equation}

Прямое применение агломеративной кластеризации ко всему множеству уникальных тегов 
требует хранения попарных расстояний и плохо масштабируется при увеличении размера словаря. 
Поэтому используется гибридная процедура кластеризации. 
На первом этапе из множества тегов выбирается подмножество $\mathcal{U}_0 \subset \mathcal{U}$, 
для которого выполняется агломеративная кластеризация по семантическому расстоянию. 
В результате формируется начальное множество кластеров

\begin{equation}
\mathcal{C}^{(0)}
=
{C_1, C_2, \ldots, C_K}.
\end{equation}

Для каждого кластера вычисляется нормированный центр $\boldsymbol{\mu}_k$.

Оставшиеся теги присоединяются к кластеру с наиболее близким центром:

\begin{equation}
c(u) = \arg\max_{k \in {1,\ldots,K}} \mathbf{e}_{u}^{\mathsf{T}}\boldsymbol{\mu}_k.
\end{equation}

Такой подход позволяет применять более точную агломеративную кластеризацию к ограниченной части тегов, 
а оставшуюся разметку обрабатывать линейно относительно числа тегов и кластеров. 
Чрезмерно крупные кластеры дополнительно разделяются, поскольку они могут объединять несколько смежных, но различных тематик. 
Ограничение размера кластера предотвращает формирование слишком общих категорий, малополезных для последующего отбора доменных документов.

Каждему кластеру $C_k$ присваивается название $v_k$. 
Оно формируется на основании множества входящих в кластер тегов и должно кратко отражать их общую семантику:

\begin{equation}
v_k = M(C_k),
\end{equation}

где $M$ --- языковая модель, используемая для именования кластеров. 
В результате формируется конечный словарь тематических классов

\begin{equation}
\mathcal{V} = {v_1,v_2,\ldots,v_K}.
\end{equation}

Каждому исходному свободному тегу сопоставляется новый тег соответствующего кластера:

\begin{equation}
\varphi(u) = v_{c(u)}.
\end{equation}

После этого исходная разметка документов преобразуется в кластеризованное пространство:

\begin{equation}
T^*(x_n) = \left\{\varphi(u) \mid u \in T(x_n) \right\}.
\end{equation}

Таким образом, синонимичные и близкие свободные теги заменяются общим тематическим классом. 
При этом один документ может одновременно относиться к нескольким темам. 
Поэтому задача обучения теггера формулируется как задача классификации по многим меткам.

Для документа $x_n$ строится бинарный вектор целевых меток

\begin{equation}
\mathbf{y}_n = (y_{n1},y_{n2},\ldots,y_{nK}) \in \{0,1\}^{K},
\end{equation}

где

\begin{equation}
y_{ni} = \mathbb{I} \left[ v_i \in T^*(x_n) \right].
\end{equation}

В качестве основы теггера используется предварительно обученный кодировщик текста. 
Для документа вычисляется контекстное представление и полученное значение пропускается через линейный слой классификатора.

\begin{equation}
\mathbf{z}_n = F(x_n),
\end{equation}

где $F$ --- модель классификатор, а $\mathbf{z}_n$ --- вектор логитов.

В отличие от многоклассовой классификации, где выходы модели нормируются совместно с помощью функции softmax, 
в задаче классификации по многим меткам вероятность каждого тега оценивается независимо:

\begin{equation}
p_{ni} = \frac{1}{1+\exp(-z_{ni})}.
\end{equation}

Итоговый выход модели имеет вид

\begin{equation}
\mathbf{p}(x_n) = (p_{n1},p_{n2},\ldots,p_{nK}),
\end{equation}

где $p_{ni}$ интерпретируется как оценка принадлежности документа $x_n$ тематическому классу $v_i$.

После обучения вероятностные оценки преобразуются в итоговый набор тегов с помощью порога $\tau$:

\begin{equation}
\widehat{y}_{i}(x)
=
\mathbb{I}
\left[
p_i(x) \geq \tau
\right],
\end{equation}

\begin{equation}
\widehat{T}(x;\tau)
=
\left\{
v_i \in \mathcal{V}
\mid
p_i(x) \geq \tau
\right\}.
\end{equation}

Снижение порога увеличивает число присваиваемых тегов и полноту разметки, но одновременно повышает число ложноположительных предсказаний. Повышение порога, напротив, увеличивает точность, но приводит к потере менее выраженных тематик. Поэтому значение $\tau$ подбирается на отложенной выборке, не использовавшейся для оптимизации параметров модели.

Для оценки качества используются точность, полнота и $F_1$-мера.
Оптимальный порог может быть определён как

\begin{equation}
\tau^{*} = \arg\max_{\tau \in \mathcal{T}} F_{1,\mathrm{micro}}(\tau),
\end{equation}

где $\mathcal{T}$ --- множество проверяемых пороговых значений.

Таким образом, предложенная методология преобразует открытую тематическую разметку языковой модели 
в фиксированное многометочное пространство. 
Крупная языковая модель применяется только на этапе построения исходной выборки и именования тематических кластеров, 
тогда как основная часть корпуса обрабатывается компактным трансформерным теггером. 
Это снижает вычислительную стоимость разметки и позволяет масштабировать её на коллекции, содержащие большое число документов.

\subsection{Определения домена документа}

Для формирования коллекции под целевой домен сначала определяется множество тегов, связанных с данным доменом. 
Такое множество может быть построено экспертно или с использованием крупной языковой модели.
Документ включается в доменную коллекцию, если его набор тегов содержит не менее заданного числа доменных тегов. 
Для документов, не удовлетворяющих этому условию, применяется дополнительная проверка на основе семантического сходства. 
Набор тегов документа преобразуется в текстовое описание и сравнивается с текстовым описанием целевого домена в эмбеддинговом пространстве. 
Если сходство превышает заданный порог, документ также включается в коллекцию.
Таким образом, отбор документов сочетает два критерия: жесткое правило по числу доменных тегов и мягкий семантический критерий.

\subsection{Оценка токсичности}

Помимо тематической принадлежности при формировании специализированных обучающих коллекций 
необходимо учитывать наличие в документах потенциально вредоносного содержания. 
В рамках настоящей работы под токсичностью понимается характеристика, 
отражающая риск присутствия в тексте агрессивной риторики, оскорблений, языка ненависти, 
описания или пропаганды опасных действий, а также другого содержания, 
нежелательного для включения в обучающую коллекцию без дополнительной проверки.

Для получения обучающей выборки используется модель-учитель, 
которая присваивает каждому документу оценку риска вредоносного содержания по порядковой шкале от 0 до 5. 
Значение 0 соответствует отсутствию признаков вредоносного содержания, 
а значение 5 --- выраженному содержанию с высоким потенциальным риском. 
Промежуточные значения отражают возрастание выраженности соответствующего признака. 

Помимо числовой оценки модель-учитель формирует краткое обоснование решения. 
Обоснования не используются непосредственно в качестве целевой переменной, 
однако позволяют контролировать качество автоматической разметки, 
выявлять неоднозначные случаи и обнаруживать систематические ошибки. 

\subsection{Оценка образовательной ценности}

Оценка образовательной ценности строится по той же схеме. 
Крупная языковая модель используется для присвоения документам оценки по дискретной шкале. 
Затем на полученной разметке обучается компактная модель, которая применяется ко всему корпусу.

Каждому документу сопоставляется значение из множества:

\begin{equation}
    e(x) \in \{0, 1, 2, 3, 4, 5\},
\end{equation}

где большие значения соответствуют большей образовательной ценности документа.

\section{Описание доменной коллекции для русского языка}

\subsection{Описание данных и параметры классификаторов}

В качестве исходного корпуса использовался раздел \texttt{rusCyrl} датасета FineWeb2. 
Документы группировались по доменам сайтов. 
Для каждого домена сохранялось не более 10 различных текстов, что позволило получить корпус объемом около 35 млн документов.
Для первичной разметки использовалась крупная языковая модель Qwen3-235B. 
Модель размечала каждый текст отдельно и присваивала ему ровно 5 тегов в свободной форме. 
Всего было размечено около 10 тыс. сайтов. 
В результате было получено около 100 тыс. уникальных исходных тегов.
Далее теги были сгруппированы с помощью эвристической процедуры кластеризации. 
После объединения близких тегов было получено 788 кластеров. 
Для каждого кластера модель Qwen выбирала репрезентативное название на основе 10 наиболее частотных тегов внутри кластера.
На полученной разметке был обучен многометочный теггер на основе RuModernBERT от VK. 
Модель возвращала вектор из 788 значений от 0 до 1. 
Тег присваивался документу, если соответствующее значение превышало выбранный порог. 
Порог подбирался на обучающей или отложенной выборке по метрикам precision, recall и F1.

\begin{table}[h]
\centering
\caption{Пример подбора порога для теггера при $k=3$}
\label{tab:tagger-threshold}
\begin{tabular}{rrrr}
\toprule
Порог & Точность & Полнота & F1 \\
\midrule
0.00 & 0.7584 & 0.5576 & 0.6427 \\
0.10 & 0.7604 & 0.5573 & 0.6432 \\
0.20 & 0.7725 & 0.5544 & 0.6455 \\
0.25 & 0.7815 & 0.5509 & 0.6463 \\
0.30 & 0.7920 & 0.5453 & 0.6459 \\
0.40 & 0.8175 & 0.5279 & 0.6416 \\
0.50 & 0.8442 & 0.4996 & 0.6277 \\
0.60 & 0.8757 & 0.4600 & 0.6032 \\
0.70 & 0.9046 & 0.4070 & 0.5615 \\
\bottomrule
\end{tabular}
\end{table}

Для дальнейших экспериментов использовался порог 0.25, обеспечивший наибольшее значение F1 в приведенной конфигурации.

\subsection{Анализ качества FineWeb2-rus по результатам разметки}

Первичный анализ распределения тегов показывает, что в корпусе присутствуют как содержательные тематические категории, 
так и значительная доля коммерческих, развлекательных и потенциально нежелательных документов. 
Среди наиболее частотных тегов встречаются туризм, онлайн-магазины, образование, строительство, дизайн интерьера, 
здоровье, производство, азартные игры, игры, акции и скидки, отели, медицина, финансы и прочее.

\subsection{Распределение документов по набору доменов}

В качестве первого примера доменной коллекции была собрана коллекция по юриспруденции. 
Для этого множество из 788 тегов было проанализировано крупной языковой моделью, которая определила теги, относящиеся к юридическому домену. 
Предварительно к юридическим было отнесено около 40 тегов.

Документ включался в юридическую коллекцию, если выполнялось хотя бы одно из двух условий:

\begin{enumerate}
    \item документ имел не менее трех тегов из юридического множества;
    \item документ содержит от 1 до 2 юридических тегов и семантическое сходство между текстовым описанием тегов документа и запросом \textit{"Тексты про юриспруденцию, законы, суды и т.д."} превышало 0.39.
\end{enumerate}

Для вычисления эмбеддингов использовалась модель \texttt{sergeyzh/BERTA}. 
Порог 0.39 был выбран эмпирически на основе предварительного просмотра результатов на нескольких текстах.
В результате было получено около 3 млн документов юридической тематики.

\subsection{Распределение фич по доменам и тегам}

\begin{table}[h]
\centering
\caption{Распределение числа тегов на документ после применения теггера}
\label{tab:tags-per-doc}
\begin{tabular}{rrr}
\toprule
Число тегов & Число документов & Доля \\
\midrule
0 & 36689 & 0.1\% \\
1 & 846632 & 2.4\% \\
2 & 3869365 & 10.9\% \\
3 & 8008849 & 22.5\% \\
4 & 9365606 & 26.3\% \\
5 & 7223097 & 20.3\% \\
6 & 3962951 & 11.1\% \\
7 & 1609156 & 4.5\% \\
8 & 516997 & 1.5\% \\
9 & 137853 & 0.4\% \\
10 & 32381 & 0.1\% \\
\bottomrule
\end{tabular}
\end{table}

\subsection{Сравнение распределений образовательной ценности для топовых моделей}

Модель Qwen3-235B присваивала текстам оценку образовательной ценности по шкале от 0 до 5. 

\section{Ограничения}

На текущем этапе качество некоторых этапов пайплайна оценивалось преимущественно автоматически или путем предварительного просмотра результатов. 
В частности, требуется дополнительная ручная валидация качества доменного отбора, качества теггера и оценок образовательной ценности. 
Также необходимо более строго обосновать выбор порогов для семантического сходства и для классификаторов.

\section{Заключение}

В статье предложен пайплайн формирования доменно-ориентированных обучающих коллекций, 
основанный на разметке крупной языковой моделью, кластеризации свободных тегов, 
обучении компактных классификаторов и последующем масштабировании разметки на большой веб-корпус. 
На примере русскоязычного корпуса FineWeb2 показано, что такой подход позволяет получать доменные коллекции 
с дополнительными признаками полезности и надежности документов.

\begin{thebibliography}{9}

\bibitem{pile}
L.~Gao et al.
\textit{The Pile: An 800GB Dataset of Diverse Text for Language Modeling}.
2020.
\url{https://arxiv.org/abs/2101.00027}

\bibitem{oscar}
J.~Abadji et al.
\textit{Towards a Cleaner Document-Oriented Multilingual Crawled Corpus}.
2022.
\url{https://aclanthology.org/2022.lrec-1.463/}

\bibitem{refinedweb}
G.~Penedo et al.
\textit{The RefinedWeb Dataset for Falcon LLM: Outperforming Curated Corpora with Web Data, and Web Data Only}.
2023.
\url{https://arxiv.org/abs/2306.01116}

\bibitem{culturax}
T.~Nguyen et al.
\textit{CulturaX: A Cleaned, Enormous, and Multilingual Dataset for Large Language Models in 167 Languages}.
2024.
\url{https://aclanthology.org/2024.lrec-main.377/}

\bibitem{fineweb}
Penedo, Guilherme, et al. "The fineweb datasets: Decanting the web for the finest text data at scale." Advances in Neural Information Processing Systems 37 (2024): 30811-30849.

\bibitem{fineweb2}
Penedo, Guilherme, et al. "FineWeb2: One Pipeline to Scale Them All--Adapting Pre-Training Data Processing to Every Language." (2025).

\end{thebibliography}

\end{document}
